%!TeX root=csp.tex
\section{Properties of strong subuniverses}\label{sec:properties}

This section is Zhuk's \S2.3. It is the \emph{interface} between the algebraic
theory and the CSP argument: \S\ref{sec:main-statements} uses the results
below and nothing else from \S\S\ref{sec:types}. Proofs are deferred to \S5 of \cite{ZhukSimplified}, which is not
reproduced here; \S\ref{sec:formalization} budgets it.

Its shape is worth naming. Four of the five groups say that the relation
$\lll$ and the types $T$ \emph{propagate}: forward and backward along
surjective homomorphisms, down to quotients, across products of a congruence,
and through subdirect relations. The fifth --- Theorem~\ref{thm:stable-intersection}
--- is the one genuinely hard statement, and it is what converts a failure of
consistency into a bridge. Everything the CSP proof does with bridges traces
back to it.

\begin{convention}\label{conv:type-quantifier}
Where a type is not specified, $T$ ranges over
$\{\TBA,\TC,\TS,\TPC,\TL,\TD\}$. Where a multi-type $\MT T$ appears,
$T$ ranges over $\{\TPC,\TL,\TD\}$.
\end{convention}

\subsection{Ubiquity}

\begin{lemma}[Ubiquity]\label{lem:ubiquity}
If $B\lll A$ and $|B|>1$, then $C\subt[A]{T}B$ for some $C$ and some
$T\in\{\TBA,\TC,\TL,\TPC\}$.
\end{lemma}

Lemma~\ref{lem:ubiquity} is the formal content of Informal
Claim~\ref{ic:existence}: every domain of size at least $2$ that has been
reached by a chain of strong reductions admits another strong reduction. It
is what guarantees the outer loop of the algorithm makes progress, and it is
the reason the algorithm terminates with singleton domains or a linear
instance.

\begin{hazard}\label{haz:ubiquity-hyp}
The hypothesis is $B\lll A$, not $B\le A$. A subalgebra reached by no chain
of strong reductions need admit nothing. In the algorithm this hypothesis is
maintained as an invariant of the reduction loop; in a formalization it should
be carried in the state, not re-derived.
\end{hazard}

\subsection{Propagation}

\begin{longlemma}[Propagation along a homomorphism]\label{lem:propagation}
Let $f\colon\mathbf A\to\mathbf A'$ be a surjective homomorphism. Then
\begin{enumerate}\tightlist
\item[\textup{(f)}] $C\lll^{A}B\Rightarrow f(C)\lll f(B)$;
\item[\textup{(b)}] $C'\lll^{A'}B'\Rightarrow f^{-1}(C')\lll f^{-1}(B')$;
\item[\textup{(ft)}] $C\subt[A]{T(\sigma)}B\lll A\Rightarrow
  \bigl(f(C)=f(B)$ or $f(C)\subt{\TS}f(B)$ or
  $f(C)\subt[A']{T}f(B)\bigr)$;
\item[\textup{(bt)}] $C'\subt[A']{T(\sigma)}B'\Rightarrow
  f^{-1}(C')\subt[A]{T(f^{-1}(\sigma))}f^{-1}(B')$;
\item[\textup{(fs)}] $T\in\{\TBA,\TC,\TS\}$ and $C\subt{T}B$
  $\Rightarrow f(C)\subte{T}f(B)$;
\item[\textup{(fm)}] $C\subte[A]{\MT T}B\lll A$ and $f(B)$ $\TS$-free
  $\Rightarrow f(C)\subte[A']{\MT T}f(B)$;
\item[\textup{(bm)}] $C'\subte[A']{\MT T}B'\lll A'\Rightarrow
  f^{-1}(C')\subte[A]{\MT T}f^{-1}(B')$.
\end{enumerate}
\end{longlemma}

\begin{hazard}[The three-way disjunction in (ft)]\label{haz:ft}
Clause (ft) is the characteristic shape of this theory and the reason it is
usable. Pushing a strong subset forward along a homomorphism can
\begin{itemize}\tightlist
\item collapse it ($f(C)=f(B)$: the reduction becomes vacuous),
\item degrade it ($f(C)\subt{\TS}f(B)$: still strong, but the type is lost),
\item or preserve it ($f(C)\subt[A']{T}f(B)$: same type).
\end{itemize}
The type $\TS$ exists to be the value of the middle case. Every downstream
proof that pushes a reduction through a quotient does a three-way case split
here, and the middle case is the one that is easy to forget; in \cite{ZhukJACM}
it is the source of the ``go forward and backward from global to local''
induction that this paper eliminates.

Note the asymmetry: the backward clauses (b), (bt), (bm) are clean
implications with no disjunction. Preimages behave; images do not.
\end{hazard}

\begin{corollary}[Propagation to a quotient]\label{cor:prop-quotient}
Let $\delta$ be a congruence on $\mathbf A$. Then
\textup{(f)} $C\lll^{A}B\Rightarrow C/\delta\lll^{A/\delta}B/\delta$;
\textup{(t)} $C\subt[A]{T(\sigma)}B\lll A\Rightarrow
 \bigl(C/\delta=B/\delta$ or $C/\delta\subt{\TS}B/\delta$ or
 $C/\delta\subt[A/\delta]{T}B/\delta\bigr)$;
\textup{(s)} for $T\in\{\TBA,\TC,\TS\}$, $C\subt{T}B\Rightarrow
 C/\delta\subte{T}B/\delta$;
\textup{(m)} $C\subte[A]{\MT T}B\lll A$ and $B/\delta$ $\TS$-free
 $\Rightarrow C/\delta\subte[A/\delta]{\MT T}B/\delta$.
\end{corollary}

\begin{corollary}[Reflection from a quotient]\label{cor:prop-from-factor}
Let $\delta$ be a congruence on $\mathbf A$ and $B,C\le\mathbf A$. Then
$C/\delta\lll^{A/\delta}B/\delta\iff C\circ\delta\lll^{A}B\circ\delta$, and
$C/\delta\subt[A/\delta]{T}B/\delta\iff C\circ\delta\subt[A]{T}B\circ\delta$.
\end{corollary}

\begin{corollary}[Multiplying by a congruence]\label{cor:prop-times-cong}
Let $\delta$ be a congruence on $\mathbf A$. Then
\textup{(f)} $C\lll^{A}B\Rightarrow C\circ\delta\lll^{A}B\circ\delta$;
\textup{(t)} $C\subt[A]{T(\sigma)}B\lll^{A}A\Rightarrow
 \bigl(C\circ\delta=B\circ\delta$ or $C\circ\delta\subt[A]{\TS}B\circ\delta$
 or $C\circ\delta\subt[A]{T}B\circ\delta\bigr)$;
\textup{(e)} if $\delta\subseteq\sigma$ and $C\subt[A]{T(\sigma)}B\lll^{A}A$
 then $C\circ\delta\subt[A]{T}B\circ\delta$.
\end{corollary}

\begin{longcorollary}[Propagation to relations]\label{cor:prop-relations}
Let $R\sdle\mathbf A_{1}\times\dots\times\mathbf A_{m}$ and $B_{i}\lll A_{i}$.
Then
\begin{enumerate}\tightlist
\item[\textup{(r)}] $R\cap(B_{1}\times\dots\times B_{m})\dlll R$;
\item[\textup{(r1)}] $\proj_{1}\bigl(R\cap(B_{1}\times\dots\times B_{m})\bigr)
  \dlll A_{1}$;
\item[\textup{(b)}] if $C_{i}\lll^{A_{i}}B_{i}$ for all $i$ then
  $R\cap(C_{1}\times\dots\times C_{m})\dlll^{R}R\cap(B_{1}\times\dots\times B_{m})$;
\item[\textup{(b1)}] under the same hypothesis,
  $\proj_{1}\bigl(R\cap\prod C_{i}\bigr)\dlll^{A_{1}}
   \proj_{1}\bigl(R\cap\prod B_{i}\bigr)$;
\item[\textup{(m)}] if $C_{i}\subte[A_{i}]{\MT T}B_{i}$ for all $i$ then
  $R\cap\prod C_{i}\dsubte[R]{\MT T}R\cap\prod B_{i}$;
\item[\textup{(m1)}] if in addition $\proj_{1}(R\cap\prod B_{i})$ is
  $\TS$-free, then $\proj_{1}(R\cap\prod C_{i})
   \dsubte[A_{1}]{\MT T}\proj_{1}(R\cap\prod B_{i})$.
\end{enumerate}
\end{longcorollary}

For the two absorption types the conclusion is stronger --- no chain, no dot on
the type, and no subdirectness hypothesis:

\begin{lemma}[{\cite[Cor.~6.1.2, 6.9.2]{ZhukStrong}}]\label{lem:ba-center-implies}
Let $R\le\mathbf A_{1}\times\dots\times\mathbf A_{m}$ with
$\proj_{1}(R)=A_{1}$, and let $C_{i}\subte{T}A_{i}$ for every $i$, where
$T\in\{\TBA,\TC\}$ --- and, when $T=\TBA$, with a \emph{single} binary term $t$
witnessing all $m$ absorptions. Then
$\proj_{1}\bigl(R\cap(C_{1}\times\dots\times C_{m})\bigr)\dsubte{T}A_{1}$,
again with the same $t$ when $T=\TBA$.
\end{lemma}

\begin{hazard}[The source's restatement is false; this is the repair]
\label{haz:ba-center-implies}
Lemma~\ref{lem:ba-center-implies} is the workhorse of case (2) of
Theorem~\ref{thm:main-inductive}, and \cite[Lemma~19]{ZhukSimplified} restates
it from \cite{ZhukStrong} \emph{dropping two hypotheses}: the subdirectness
$\proj_{1}(R)=A_{1}$, and --- in the $\TBA$ case --- that the $C_{i}$ absorb
with a common term. Both are present in the cited source
(\cite[Cor.~6.1.2]{ZhukStrong} reads ``Suppose $R\le A_{1}\times\dots\times
A_{n}$, $\proj_{1}(R)=A_{1}$, $B_{i}\le_{\TBA(t)}A_{i}$ for every $i$''), and
without the first the statement is false.

\emph{Counterexample to the printed form.} Take $R=\{(0,0)\}\le\mathbf Z_{p}
\times\mathbf Z_{p}$ --- a subuniverse, by idempotence --- and $C_{i}=A_{i}$,
which is allowed because $\subte{T}$ admits equality. Then
$R\cap(C_{1}\times C_{2})=R$ and $\proj_{1}(R)=\{0\}$, which is neither empty
nor all of $\mathbf Z_{p}$, so the conclusion asserts $\{0\}\subt{\TBA}
\mathbf Z_{p}$ --- contradicting Lemma~\ref{lem:no-abs-in-linear}, the paper's
own Lemma~29. The failure is exactly the missing subdirectness: $\proj_{1}(R)
=\{0\}\neq\mathbf Z_{p}$.

The repair is to restore both hypotheses, as above. A formalization must also
check that they hold at each of the call sites --- in particular in the proof of
Theorem~\ref{thm:main-inductive}(2), where the lemma is applied to the relation
computing $E_{2}$ from $E_{1}$. The common-term condition is the reason
\cite{ZhukStrong} writes $\le_{\TBA(t)}$ with the term in the notation, and
carrying $t$ in the type is the cheapest way to keep it.

This is the single highest-value import to formalize first: small,
self-contained, used everywhere, and adjacent to what the predecessor project
already proves.
\end{hazard}

\subsection{Intersections}

\begin{lemma}[Intersecting with a strong subset]\label{lem:intersect-all}
If $B\lll A$ and $D\lll A$ then \textup{(i)} $B\cap D\dlll A$, and
\textup{(t)} $C\subt[A]{T(\sigma)}B\Rightarrow C\cap D\dsubte[A]{T(\sigma)}B\cap D$.
\end{lemma}

The next theorem is the heart of \S\ref{sec:properties}. Read it as follows:
if a family of strong reductions is \emph{jointly} inconsistent but
\emph{minimally} so --- dropping any one of them restores consistency --- then
the types must agree, and in the linear case one gets bridges between the
witnessing congruences. Bridges are produced nowhere else.

\begin{longtheorem}[Stable intersection]\label{thm:stable-intersection}
Suppose
\begin{enumerate}\romanlabels\tightlist
\item $C_{i}\subt[A]{T_{i}(\sigma_{i})}B_{i}\lll A$ with
  $T_{i}\in\{\TBA,\TC,\TS,\TL,\TPC\}$, for $i\in[m]$, $m\ge 2$;
\item $\bigcap_{i\in[m]}C_{i}=\varnothing$;
\item $B_{j}\cap\bigcap_{i\neq j}C_{i}\neq\varnothing$ for every $j\in[m]$.
\end{enumerate}
Then one of:
\begin{enumerate}\tightlist
\item[\textup{(ba)}] $T_{1}=\dots=T_{m}=\TBA$;
\item[\textup{(l)}] $T_{1}=\dots=T_{m}=\TL$, and for all $k,\ell\in[m]$ there
  is a bridge $\delta$ from $\sigma_{k}$ to $\sigma_{\ell}$ with
  $\bcol{\delta}=\sigma_{k}\circ\sigma_{\ell}$;
\item[\textup{(c)}] $m=2$ and $T_{1}=T_{2}=\TC$;
\item[\textup{(pc)}] $m=2$, $T_{1}=T_{2}=\TPC$ and $\sigma_{1}=\sigma_{2}$.
\end{enumerate}
\end{longtheorem}

\begin{longcorollary}[Stable intersection, relational form]\label{cor:stable-intersection}
Suppose $R\sdle\mathbf A_{1}\times\dots\times\mathbf A_{m}$,
$C_{i}\subt[A_{i}]{T_{i}(\sigma_{i})}B_{i}\lll A_{i}$ with
$T_{i}\in\{\TBA,\TC,\TS,\TL,\TPC\}$ and $m\ge 2$,
$R\cap(C_{1}\times\dots\times C_{m})=\varnothing$, and
$R\cap(C_{1}\times\dots\times B_{j}\times\dots\times C_{m})\neq\varnothing$ for
every $j$. Then one of:
\textup{(ba)} all $T_{i}=\TBA$;
\textup{(l)} all $T_{i}=\TL$ and for all $k,\ell$ a bridge $\delta$ from
$\sigma_{k}$ to $\sigma_{\ell}$ with
$\bcol{\delta}=\sigma_{k}\circ\proj_{k,\ell}(R)\circ\sigma_{\ell}$;
\textup{(c)} $m=2$ and $T_{1}=T_{2}=\TC$;
\textup{(pc)} $m=2$, $T_{1}=T_{2}=\TPC$,
$\mathbf A_{1}/\sigma_{1}\cong\mathbf A_{2}/\sigma_{2}$, and
$\{(a/\sigma_{1},b/\sigma_{2}):(a,b)\in R\}$ is bijective.
\end{longcorollary}

\begin{remark}\label{rem:duplicate-coordinate}
Several applications need more than one restriction on a single coordinate of
$R$. The statement is kept simple; to apply it, duplicate the coordinate (form
$R'=\{(\bar a,a_{k}):\bar a\in R\}$) and restrict the copies separately.
\end{remark}

\begin{hazard}[Minimality is a hypothesis, not a construction]\label{haz:minimality}
Hypotheses (ii)--(iii) of Theorem~\ref{thm:stable-intersection} say the family
is \emph{critically} inconsistent. In applications one starts with an
inconsistent family and shrinks it to a critical subfamily. That shrinking step
is left implicit in the paper (``let $\mathcal M$ be a minimal set of children
of $z$ we need to restrict''). It is easy --- a minimal subset of a finite set
with a monotone property --- but it must be stated as a lemma, because it is
performed at least four times and each time the property being minimised is
slightly different. State once: \emph{if $P$ is a monotone predicate on subsets
of a finite set $X$ and $P(X)$ holds, there is $M\subseteq X$ minimal with
$P(M)$}, and instantiate.
\end{hazard}

\subsection{Multi-types}

\begin{lemma}\label{lem:multitype-chain}
If $C\subte[A]{\MT T}B$ then $C\subt[A]{T}\dots\subt[A]{T}B$ and
$C\lll^{A}B$.
\end{lemma}

That is: an intersection of type-$T$ subsets, though not itself obtained by a
single reduction of type $T$, is reachable by a \emph{chain} of type-$T$
reductions. This is what makes $\MT T$ compatible with $\lll$, and it is used
whenever the algorithm reduces several variables at once.

\begin{lemma}[Preserving linkedness]\label{lem:preserve-linked}
Let $R\sdle\mathbf A_{1}\times\mathbf A_{2}$,
$C_{i}\subte[A_{i}]{\TMD}B_{i}\lll A_{i}$ for $i=1,2$, let $S$ be the
rectangular closure of $R$, and suppose $R\cap(B_{1}\times B_{2})\neq
\varnothing$ and $S\cap(C_{1}\times C_{2})\neq\varnothing$. Then
$R\cap(C_{1}\times C_{2})\neq\varnothing$.
\end{lemma}

\begin{lemma}[Maximal multi-type extension]\label{lem:maximal-mult}
Let $C_{1}\subt[A]{\MT T}B_{1}\lll A$, $B_{2}\lll A$,
$C_{1}\cap B_{2}=\varnothing$, $B_{1}\cap B_{2}\neq\varnothing$, and let
$\sigma$ be a maximal congruence on $\mathbf A$ with
$(C_{1}\circ\sigma)\cap B_{2}=\varnothing$. Then
$\sigma=\omega_{1}\cap\dots\cap\omega_{s}$ where each $\omega_{i}$ is a
congruence of type $T$ on $\mathbf A$ with $\cover{\omega_{i}}\supseteq
B_{1}^{2}$.
\end{lemma}

\subsection{Auxiliary imports}

The following are used but not proved in \cite{ZhukSimplified}.

\begin{lemma}[{\cite[Lem.~4.7]{MarotiMcKenzie}}]\label{lem:special-wnu}
If $w$ is an idempotent WNU operation of arity $n$ on a finite set $A$, then
$\Clo(w)$ contains a special idempotent WNU operation of arity $n^{n!}$.
\end{lemma}

\begin{lemma}[{\cite[Cor.~8.17.1]{ZhukJACM}}]\label{lem:building-perfect}
If $\sigma$ is an irreducible congruence on $\mathbf A\in\Vn$ and $\delta$ is a
bridge from $\sigma$ to $\sigma$ with $\bcol{\delta}=A^{2}$, then $\sigma$ is a
perfect linear congruence.
\end{lemma}

\begin{lemma}\label{lem:perfect-from-linked}
If $\sigma$ is an irreducible congruence on $\mathbf A\in\Vn$ and $\delta$ is a
bridge from $\sigma$ to $\sigma$ with $\bcol{\delta}$ linked, then $\sigma$ is
a perfect linear congruence.
\end{lemma}

\begin{proof}
Let $\delta^{-1}(y_{1},y_{2},x_{1},x_{2})=\delta(x_{1},x_{2},y_{1},y_{2})$,
which is again a bridge from $\sigma$ to $\sigma$. Since $\bcol{\delta}$ is
linked, $(\bcol{\delta}\circ\bcol{\delta^{-1}})^{N}=A^{2}$ for $N$ large.
Composing $2N$ bridges $\delta,\delta^{-1},\dots$ by
Lemma~\ref{lem:bridge-comp} gives a bridge $\delta'$ from $\sigma$ to $\sigma$
with $\bcol{\delta'}=A^{2}$; apply Lemma~\ref{lem:building-perfect}.
\end{proof}

\begin{hazard}[``for sufficiently large $N$'']\label{haz:linked-power}
The step $(\bcol{\delta}\circ\bcol{\delta^{-1}})^{N}=A^{2}$ deserves a lemma of
its own: if a reflexive symmetric relation $\rho$ on a finite set $A$ has
connected graph, then $\rho^{N}=A^{2}$ for $N\ge|A|$. Reflexivity is what makes
the powers increase monotonically; $\bcol{\delta}$ is reflexive because
$(a,a,a,a)\in\delta$ is not assumed --- so the reflexivity must come from
$\delta$ being a bridge from $\sigma$ to $\sigma$ with $\sigma$ reflexive, via
$\proj_{1,2}(\delta)\supsetneq\sigma$. Check this: it is the sort of step where
the informal proof is right for a reason it does not give.
\end{hazard}

\begin{lemma}\label{lem:no-abs-in-linear}
If $B\le\mathbf Z_{p}\times\dots\times\mathbf Z_{p}$ then there is no
$C\subt{T}B$ with $T\in\{\TBA,\TC\}$.
\end{lemma}

\begin{proof}
It suffices to check that for every term $\tau$, if the last variable of
$\tau^{\mathbf Z_{p}}(a,\dots,a,x)$ is not dummy then it takes all values;
hence no absorbing subuniverse exists.
\end{proof}
